% SVM_CIFAR10_Report.tex
% XeLaTeX report for final assignment: SVM and MLP for CIFAR-10
% Compile with: xelatex SVM_CIFAR10_Report.tex

\documentclass[11pt,a4paper]{article}

% XeLaTeX packages for Greek support
\usepackage{fontspec}
\usepackage{polyglossia}
\setmainlanguage{greek}
\setotherlanguage{english}
\setmainfont{Liberation Serif}
\setsansfont{Liberation Sans}
\setmonofont{Liberation Mono}
\newfontfamily\greekfonttt{Liberation Mono}

% Additional packages
\usepackage{graphicx}
\usepackage{geometry}
\geometry{margin=2.5cm}
\usepackage{fancyhdr}
\usepackage{booktabs}
\usepackage{amsmath}
\usepackage{hyperref}
\usepackage{float}
\hypersetup{
    colorlinks=true,
    linkcolor=blue,
    filecolor=magenta,      
    urlcolor=cyan,
}

% Custom page style
\pagestyle{fancy}
\fancyhf{}
\fancyfoot[R]{\thepage}
\renewcommand{\headrulewidth}{0pt}
\renewcommand{\footrulewidth}{0pt}

\begin{document}

% ======================== TITLE PAGE ========================
\begin{center}
  % AUTH Logo
  \IfFileExists{AUTH_logo.png}{%
    \vspace*{-1.5cm}
    \includegraphics[width=0.85\textwidth]{AUTH_logo.png}\\[1.5em]
  }{%
    \vspace*{-1.0cm}
    {\Large \textit{AUTH logo}}\\[1.5em]
  }

  % Student Information
  {\fontsize{22pt}{26pt}\selectfont\textbf{Ευάγγελος Μόσχου}}\\[0.5em]
  {\fontsize{16pt}{20pt}\selectfont\textbf{}}\\[1.2em]
  
  {\fontsize{14pt}{18pt}\selectfont Τμήμα Ηλεκτρολόγων Μηχανικών και Μηχανικών Υπολογιστών}\\[0.3em]
  {\fontsize{14pt}{18pt}\selectfont Αριστοτέλειο Πανεπιστήμιο Θεσσαλονίκης}\\[0.5em]
  {\fontsize{14pt}{18pt}\selectfont Νοέμβριος 2025}\\[2.0em]
\end{center}

\vspace*{0.35\textheight}

\begin{center}
  {\fontsize{20pt}{24pt}\selectfont\textbf{%
    Υποχρεωτική Εργασία 2\\[0.5em]
    Νευρωνικά Δίκτυα και Βαθιά Μάθηση}}\\[1.5em]
  {\fontsize{18pt}{22pt}\selectfont%
    Ταξινόμηση εικόνων CIFAR-10\\[0.3em]
    με χρήση SVM και MLP}
\end{center}

\newpage

% ======================== ΠΕΡΙΕΧΟΜΕΝΑ ========================
\tableofcontents
\newpage

% ======================== 1. ΕΙΣΑΓΩΓΗ ========================
\section{Εισαγωγή}

\subsection{Περιγραφή του προβλήματος}
Η παρούσα εργασία αφορά την υλοποίηση και αξιολόγηση αλγορίθμων μηχανικής μάθησης για την ταξινόμηση εικόνων από το σύνολο δεδομένων \textenglish{CIFAR-10}. Συγκεκριμένα, εξετάζονται οι αλγόριθμοι \textenglish{Support Vector Machines (SVM)} και \textenglish{Multi-Layer Perceptron (MLP)} με χρήση \textenglish{Hinge Loss}, καθώς και η σύγκρισή τους με \textenglish{K-Nearest Neighbors (KNN)} και \textenglish{Nearest Class Centroid (NCC)}.

\subsection{Το σύνολο δεδομένων CIFAR-10}
Το \textenglish{CIFAR-10} αποτελείται από 60.000 έγχρωμες εικόνες διαστάσεων $32 \times 32 \times 3$, κατανεμημένες σε 10 κλάσεις. Για την παρούσα εργασία, τα δεδομένα υποβλήθηκαν σε προεπεξεργασία που περιλαμβάνει κανονικοποίηση και μείωση διαστάσεων με \textenglish{Principal Component Analysis (PCA)}.

\subsection{Στόχος της εργασίας}
Ο στόχος είναι η συγκριτική μελέτη των παραπάνω αλγορίθμων ως προς την ακρίβεια ταξινόμησης και τον χρόνο εκπαίδευσης, καθώς και η διερεύνηση της επίδρασης των υπερπαραμέτρων (π.χ. πυρήνας SVM, πλήθος νευρώνων MLP).

\newpage

% ======================== 2. ΜΕΘΟΔΟΛΟΓΙΑ ========================
\section{Μεθοδολογία}

\subsection{Προεπεξεργασία Δεδομένων}
Οι εικόνες μετατράπηκαν σε διανύσματα μεγέθους 3072 ($32 \times 32 \times 3$) και κανονικοποιήθηκαν στο διάστημα $[0, 1]$. Στη συνέχεια εφαρμόστηκε \textenglish{PCA} διατηρώντας το 90\% της διακύμανσης, μειώνοντας σημαντικά τις διαστάσεις του χώρου χαρακτηριστικών.

\subsection{Support Vector Machines (SVM)}
Δοκιμάστηκαν διάφοροι πυρήνες (\textenglish{kernels}):
\begin{itemize}
    \item \textenglish{Linear}
    \item \textenglish{RBF (Radial Basis Function)}
    \item \textenglish{Polynomial}
    \item \textenglish{Sigmoid}
\end{itemize}
Για κάθε πυρήνα πραγματοποιήθηκε αναζήτηση πλέγματος (\textenglish{grid search}) για τις παραμέτρους $C$ και $\gamma$.

\subsection{Multi-Layer Perceptron (MLP)}
Υλοποιήθηκε ένα δίκτυο MLP με ένα κρυφό επίπεδο, το οποίο βελτιστοποιήθηκε σημαντικά σε σχέση με την αρχική απλή υλοποίηση.
\begin{itemize}
    \item \textbf{Αρχιτεκτονική}: Input $\rightarrow$ Dense (512) $\rightarrow$ Batch Normalization $\rightarrow$ ReLU $\rightarrow$ Dropout (0.3) $\rightarrow$ Output (10)
    \item \textbf{Bελτιστοποίηση}: Χρήση του \textbf{AdamW} (Adam με Weight Decay) για καλύτερη κανονικοποίηση.
    \item \textbf{Techniques}: Εφαρμόστηκε \textbf{Early Stopping} και μείωση ρυθμού μάθησης (\textbf{ReduceLROnPlateau}) όταν η απώλεια σταθεροποιείται.
    \item \textbf{Συνάρτηση Κόστους}: \textenglish{Hinge Loss} (όπως ζητήθηκε).
\end{itemize}

\subsection{Άλλοι Ταξινομητές}
Για λόγους σύγκρισης, χρησιμοποιήθηκαν επίσης:
\begin{itemize}
    \item \textenglish{K-Nearest Neighbors} (1-NN και 3-NN)
    \item \textenglish{Nearest Class Centroid (NCC)}
\end{itemize}

\newpage

% ======================== 3. ΑΠΟΤΕΛΕΣΜΑΤΑ ========================
\section{Αποτελέσματα}

\subsection{Ανάλυση PCA}
Για τα Deep Features (MobileNetV2), η εφαρμογή του PCA μείωσε τις διαστάσεις από 1280 σε \textbf{475 συνιστώσες}, διατηρώντας το \textbf{90\%} της αρχικής διακύμανσης.

\begin{figure}[H]
\centering
\includegraphics[width=0.7\textwidth]{../results/figures/pca_variance.png}
\caption{Σωρευτική εξηγούμενη διακύμανση (Deep Features)}
\label{fig:pca}
\end{figure}

\subsection{Σύγκριση Ταξινομητών}
Ο Πίνακας \ref{tab:results} παρουσιάζει τα συγκεντρωτικά αποτελέσματα για όλους τους ταξινομητές.

\begin{table}[H]
\centering
\caption{Συγκεντρωτικά αποτελέσματα ακρίβειας και χρόνου εκπαίδευσης}
\label{tab:results}
\resizebox{\textwidth}{!}{%
\begin{tabular}{@{}llcc@{}}
\toprule
\textbf{Ταξινομητής} & \textbf{Διαμόρφωση} & \textbf{Ακρίβεια (Test \%)} & \textbf{Χρόνος Εκπ. (s)} \\ 
\midrule
\textbf{SVM (Deep Features)} & \textbf{Linear (MobileNetV2 + PCA 90\%)} & \textbf{85.94\%} & \textbf{$\sim$40} \\
\midrule
SVM & RBF ($C=10, \gamma=0.01$, Raw Pixels) & 56.18\% & 10.22 \\
MLP (Hinge) & 512 νευρώνες & 53.47\% & 14.69 \\
KNN & $k=1$ & 38.58\% & 0.004 \\
NCC & Centroids & 27.66\% & 0.039 \\
\bottomrule
\end{tabular}%
}
\end{table}

\subsection{Οπτικοποίηση Σύγκρισης}
Το παρακάτω διάγραμμα απεικονίζει τη συγκριτική απόδοση. Η χρήση \textenglish{Deep Features} (MobileNetV2) πέτυχε την υψηλότερη ακρίβεια \textbf{85.94\%}, ξεπερνώντας κατά πολύ τους παραδοσιακούς ταξινομητές.

\begin{figure}[H]
\centering
\includegraphics[width=1.0\textwidth]{../results/figures/classifier_comparison.png}
\caption{Συγκριτικό ραβδόγραμμα ακρίβειας}
\label{fig:comparison}
\end{figure}

\subsection{Πίνακας Σύγχυσης (Confusion Matrix)}
Ο πίνακας σύγχυσης του καλύτερου μοντέλου (Deep Features + Linear SVM):

\begin{figure}[H]
\centering
\includegraphics[width=0.8\textwidth]{../results/figures/best_model_confusion_matrix.png}
\caption{Πίνακας Σύγχυσης για το Deep Features SVM}
\label{fig:cm}
\end{figure}

\subsection{Παραδείγματα Ταξινόμησης}
Ακολουθούν παραδείγματα από το σύνολο ελέγχου όπου το μοντέλο SVM (Linear) προέβλεψε σωστά και λανθασμένα.

\begin{figure}[H]
\centering
\includegraphics[width=1.0\textwidth]{../results/figures/correct_examples.png}
\caption{Παραδείγματα σωστής κατηγοριοποίησης}
\label{fig:correct}
\end{figure}

\begin{figure}[H]
\centering
\includegraphics[width=1.0\textwidth]{../results/figures/misclassified_examples.png}
\caption{Παραδείγματα λανθασμένης κατηγοριοποίησης}
\label{fig:incorrect}
\end{figure}

\newpage

% ======================== 4. ΣΥΜΠΕΡΑΣΜΑΤΑ ========================
\section{Συμπεράσματα}
Από την πειραματική διαδικασία προκύπτουν τα εξής συμπεράσματα:

\begin{enumerate}
    \item \textbf{Βέλτιστο Μοντέλο}: Ο \textbf{SVM} με διατήρηση πληροφορίας 95\% (PCA), $C=10$ και $\gamma=0.005$ πέτυχε την κορυφαία ακρίβεια \textbf{56.94\%}. Η αύξηση της διακύμανσης του PCA παρείχε πλουσιότερα χαρακτηριστικά, επιτρέποντας στον SVM να ξεπεράσει οριακά το βελτιστοποιημένο MLP (\textbf{56.45\%}).
     
    \item \textbf{MLP Optimization}: Η προσθήκη των \textbf{AdamW}, \textbf{Batch Normalization} και \textbf{Dropout} βελτίωσε σημαντικά την απόδοση του MLP (από 53.5\% σε 56.5\%), καθιστώντας το εξαιρετικά ανταγωνιστικό με τον SVM.
    
    \item \textbf{Επιτάχυνση με GPU}: Η χρήση της βιβλιοθήκης \texttt{cuML} και ειδικά του \texttt{LinearSVC} μείωσε δραματικά τους χρόνους εκπαίδευσης (από λεπτά σε κλάσματα δευτερολέπτου για το Linear SVM), επιτρέποντας την ταχεία διερεύνηση υπερπαραμέτρων.
    
    \item \textbf{MLP με Hinge Loss}: Το δίκτυο MLP απέδωσε αρκετά καλά (\textbf{53.47\%}), πλησιάζοντας τον SVM. Η χρήση Hinge Loss αντί για Cross-Entropy λειτούργησε σωστά, προσομοιώνοντας τη συμπεριφορά "Maximum Margin" του SVM σε περιβάλλον νευρωνικών δικτύων.
    
    \item \textbf{Επίδραση PCA}: Η μείωση διαστάσεων (90\% διακύμανση) βοήθησε στην ταχύτητα, αλλά πιθανώς περιόρισε τη μέγιστη δυνατή ακρίβεια των γραμμικών μοντέλων, καθώς χάθηκε λεπτομερής πληροφορία των εικόνων.
\end{enumerate}

\end{document}
